\section{Related Work}
\label{sec:related_work}

The literature on model ensembling and multi-task learning has evolved rapidly, progressing from static parameter aggregation to dynamic, input-dependent activation-space ensembling. This section situates our methodological audit within these key domains and details the underlying architectures we scrutinize.

\subsection{Static Parameter-Space Model Merging}
Static parameter-space merging seeks to combine multiple task-specific neural network weights into a single, unified model without additional training. Early approaches relied on simple weight averaging, which often degrades performance due to parameter interference and the non-convex nature of the loss landscape \citep{touvron23llama, dosovitskiy20vit}. To address this, weight alignment and permutation-matching techniques (e.g., RE-basin, OT-fusion) align neuron representations before averaging. 

Other techniques, such as Fisher Merging and Task Arithmetic, scale task-specific parameter updates based on empirical Fisher information or task vector subtraction \citep{hu2021lora}. While parameter-space merging enjoys zero-shot inference efficiency, it remains a \emph{static} compromise: once the models are merged, the ensembling weights are fixed. When facing highly heterogeneous, non-stationary streams of data, static models are fundamentally bounded by parameter interference and cannot adapt to sample-level variations.

\subsection{Stateless Activation-Space Blending}
To overcome the limitations of static merging, researchers have turned to dynamic, activation-space blending. This paradigm is inspired by Mixture-of-Experts (MoE) architectures, but operates at serving-time over complete pre-trained adapters rather than routing token-by-token within sparse layers. Instead of running multiple sequential passes through $K$ adapters, these methods pass activations through the base network and route them through parallel adapter paths, blending their output activations dynamically at each layer.

SABLE (Sample-wise Activation Blending of Low-Rank Experts) \citep{liang23sable} represents a prominent stateless approach. It uses early-layer activation features to perform a training-free, sample-wise nearest-centroid routing. Specifically, SABLE computes the cosine similarity between the current sample representation and pre-computed task centroids to determine the blending coefficients $\boldsymbol{\alpha}_b$ at each layer. 

Similarly, PAC-ZCA \citep{zhang24spszca} introduces PAC-Bayesian bound minimization to optimize routing temperatures, and whitening (ZCA) to handle representation anisotropy. However, a major claim in these works is that classical parametric linear routers (trained via gradient descent on small calibration sets) catastrophically collapse or severely overfit, necessitating training-free, heuristic-driven geometric projection methods.

\subsection{Stateful Continuous-Time Routing}
Recognizing that layer-wise independent routing can cause severe routing weight oscillations ("jitter") and representational drift across deep layers, recent work has proposed stateful, continuous-time routing frameworks. ChemMerge \citep{chemmerge25} is the prime representative of this class. It models the routing trajectory as a continuous-time dynamical system governed by first-order chemical reaction kinetics. By formulating the layer-wise ensembling weights via concentrations $C_k(t)$ of chemical species evolving via an ordinary differential equation (ODE) \citep{chen2018neuralode}, ChemMerge applies a temporal low-pass filter to the trajectory:
\begin{equation}
    \label{eq:chem_kinetics}
    \frac{dC_k(t)}{dt} = R_k(t) (1 - C_k(t)) - K_{\text{decay}} C_k(t)
\end{equation}
where $t$ corresponds to the layer index scaled continuously, $K_{\text{decay}}$ is the chemical reaction decay rate, and $R_k(t) = \text{Softmax}_k(\text{sims}_t / \tau)$ represents the reaction rate based on nearest-centroid similarity projections against task prototypes. The final layer-wise ensembling weights $\boldsymbol{\alpha}(t)$ are obtained by normalizing the concentrations to sum to unity: $\alpha_k(t) = C_k(t) / \sum_{j=1}^K C_j(t)$. 

This formulation successfully dampens layer-to-layer routing weight jitter, and prior work claims that this "kinetic smoothing" preserves intermediate feature quality and improves overall accuracy. However, this stateful inertia is highly complex to implement, requires numerical ODE solvers (such as discretized Euler steps) at serving time, and introduces significant computational overhead.

Furthermore, we note a critical numerical stability issue in standard implementations of ChemMerge: the discretized Euler solver uses an extremely large step size (typically $\Delta t = 1.5$). In classical numerical analysis, such a large step size is highly unstable and prone to overshoot (e.g., resulting in chemical concentrations $C_k(t) > 1.0$ or $C_k(t) < 0.0$). To prevent numerical divergence, the practical implementation employs a "numerical hack"---hard-clamping the concentrations to the range $[0.0, 1.0]$ after each step. Acknowledging this numerical clamping is crucial, as it demystifies the continuous kinetics metaphor and exposes it as a hand-crafted, discrete heuristic that relies on ad-hoc constraints to maintain stability.

\subsection{The Necessity of Methodological Audits}
Our work is motivated by a long-standing tradition of methodological skepticism in machine learning. In fields ranging from recommender systems to graph neural networks, researchers have repeatedly exposed that highly complex, mathematically elaborate SOTA architectures are frequently outperformed by properly tuned classical baselines (e.g., simple linear layers or classical matrix factorization) when evaluated under rigorous, controlled conditions \citep{langley00, mitchell80, kearns89}. 

In the model merging literature, the "collapse" of classical linear routers has been widely reported. However, we observe a critical confounding variable: these classical baselines are almost universally initialized with random weights and trained on tiny calibration splits ($N \le 128$) without proper regularization (such as L2 weight decay or maximum-entropy initialization). This under-regularized setup is a textbook recipe for overfitting. By applying a rigorous methodological audit, we systematically isolate this confounding factor to uncover whether complex chemical or physical metaphors are truly necessary, or if they are simply compensating for poorly optimized baselines.
